%! The Fundamental Theorem of Linear Algebra — Unit 3, Lesson 3.6 — Guided Notes
\documentclass[10pt]{article}
\usepackage{linalg-article}
\usepackage{linalg-boxes}
\newcommand{\TallMath}[1]{$\displaystyle #1\rule[-1.4em]{0pt}{3.2em}$}
\newcommand{\vv}{\mathbf{v}}
\newcommand{\ww}{\mathbf{w}}
\newcommand{\xx}{\mathbf{x}}
\newcommand{\yy}{\mathbf{y}}
\newcommand{\aaa}{\mathbf{a}}
\newcommand{\bb}{\mathbf{b}}
\newcommand{\svec}{\mathbf{s}}
\newcommand{\zero}{\mathbf{0}}
\newcommand{\T}{\mathsf{T}}

\begin{document}

\pageheader{Unit 3, Lesson 3.6}{Guided Notes}
\namedateperiod

\vspace{0.06in}

\begin{objectivebox}
\textbf{By the end of this lesson, I will be able to\dots}
\begin{itemize}\setlength{\itemsep}{1pt}
  \item draw the \emph{big picture} --- two rooms, four subspaces, four dimensions ($r$, $n-r$, $r$, $m-r$);
  \item state \emph{Part~1} (the dimensions) and \emph{Part~2} (the row space $\perp$ nullspace, the column space $\perp$ left nullspace);
  \item explain why $A\xx=\zero$ makes $\xx$ perpendicular to every row, and \emph{check} a right angle with dot products.
\end{itemize}
\end{objectivebox}

\vspace{0.04in}

\begin{vocabbox}
\small
Fill in each term as we build it together.
\vspace{2pt}
\termblanklong{Fundamental Theorem, Part 1}
\termblanklong{Fundamental Theorem, Part 2}
\termblanklong{Orthogonal complement}
\termblanklong{The big picture}
\end{vocabbox}

\begin{hookbox}
Across the unit we met the four subspaces one at a time. Today we put them in \textbf{one picture} --- and
find a surprise.
\begin{itemize}\setlength{\itemsep}{1pt}
  \item When you compute $A\xx$, what is entry $i$ (in terms of a row of $A$ and $\xx$)? \writeline
  \item So if $A\xx=\zero$, what is true about \emph{every} row of $A$ and $\xx$? \writeline
\end{itemize}
\end{hookbox}

\begin{notesbox}{1. The Big Picture (Part 1: the dimensions)}
An $m\times n$ matrix $A$ of rank $r$ connects two rooms, and splits each into two fundamental subspaces:
\begin{center}
\begin{tikzpicture}[scale=0.92, every node/.style={font=\scriptsize}]
  % input room
  \draw[burgundy, very thick, rounded corners] (0,0) rectangle (3.4,4);
  \node[burgundy, font=\bfseries\footnotesize] at (1.7,4.32) {input room $\mathbb{R}^n$};
  \draw[burgundy, thick] (0,2) -- (3.4,2);
  \node[burgundy, font=\bfseries\scriptsize] at (1.7,3.15) {row space $C(A^{\T})$};
  \node[burgundy] at (1.7,2.6) {$\dim = r$};
  \node[burgundy, font=\bfseries\scriptsize] at (1.7,1.35) {nullspace $N(A)$};
  \node[burgundy] at (1.7,0.8) {$\dim = n-r$};
  \node[fill=white, text=burgundy, inner sep=1.2pt] at (1.7,2) {$\perp$};
  % output room
  \draw[royalblue, very thick, rounded corners] (6.1,0) rectangle (9.5,4);
  \node[royalblue, font=\bfseries\footnotesize] at (7.8,4.32) {output room $\mathbb{R}^m$};
  \draw[royalblue, thick] (6.1,2) -- (9.5,2);
  \node[royalblue, font=\bfseries\scriptsize] at (7.8,3.15) {column space $C(A)$};
  \node[royalblue] at (7.8,2.6) {$\dim = r$};
  \node[royalblue, font=\bfseries\scriptsize] at (7.8,1.35) {left nullspace $N(A^{\T})$};
  \node[royalblue] at (7.8,0.8) {$\dim = m-r$};
  \node[fill=white, text=royalblue, inner sep=1.2pt] at (7.8,2) {$\perp$};
  % arrow A
  \draw[->, ultra thick, charcoal] (3.6,2) -- (5.9,2);
  \node[charcoal, font=\bfseries\small] at (4.75,2.35) {$A$};
\end{tikzpicture}
\end{center}
Two subspaces live in the \textbf{input room} $\mathbb{R}^n$ and two in the \textbf{output room}
$\mathbb{R}^m$. \textbf{Part~1} of the Fundamental Theorem gives every dimension from the one number $r$:
\[
  \dim C(A)=\blank{0.7cm},\quad \dim C(A^{\T})=\blank{0.7cm},\quad
  \dim N(A)=\blank{1.3cm},\quad \dim N(A^{\T})=\blank{1.3cm},
\]
and the two counting rules $r+(n-r)=n$ and $r+(m-r)=m$ say each room is filled. (All of this is Lesson~3.5 ---
now drawn as \emph{one} diagram.) The new idea is the little $\perp$ in each room.
\end{notesbox}

\begin{notesbox}{2. The Right-Angle Surprise (Part 2)}
Why is the nullspace \emph{perpendicular} to the row space? One line. Compute $A\xx$ \textbf{row by row} ---
entry $i$ is (row$_i$ of $A$)$\,\cdot\,\xx$. So
\[
  A\xx=\zero \quad\Longleftrightarrow\quad \text{every row}_i\cdot\xx=\blank{0.7cm}.
\]
By Lesson~1.2, a dot product of $0$ means the two vectors are \blank{2.6cm}. So a nullspace vector $\xx$ is
perpendicular to \emph{every} row of $A$ --- and therefore to every combination of the rows, i.e.\ to the
whole \blank{2.2cm}. In symbols:
\[
  N(A)\ \perp\ C(A^{\T}) \quad\text{in the input room }\mathbb{R}^n.
\]
Nothing new was needed --- just ``$A\xx$ is a stack of dot products'' plus ``dot $=0$ means perpendicular.''
\end{notesbox}

\begin{notesbox}{3. The Other Pair, and Orthogonal Complements}
Run the \emph{same} argument on $A^{\T}$. A left-nullspace vector $\yy$ satisfies $A^{\T}\yy=\zero$, so $\yy$ is
perpendicular to every row of $A^{\T}$ --- that is, every \blank{1.8cm} of $A$. Hence
\[
  N(A^{\T})\ \perp\ C(A) \quad\text{in the output room }\mathbb{R}^m.
\]
In each room the two subspaces do \emph{two} things at once: they are perpendicular \textbf{and} their
dimensions add to fill the room ($r+(n-r)=n$; $r+(m-r)=m$). A pair like that meets only at $\zero$ and
together spans the room --- each is ``everything perpendicular to the other.'' We call them \textbf{orthogonal
complements}:
\begin{center}\small
In $\mathbb{R}^n$: row space and nullspace are orthogonal complements.\\
In $\mathbb{R}^m$: column space and left nullspace are orthogonal complements.
\end{center}
\end{notesbox}

\begin{notesbox}{4. The Whole Theorem, on One Matrix}
Take the unit's running matrix $A=\begin{bmatrix} 1 & 1 & 2 \\ 2 & 1 & 3 \\ 3 & 2 & 5 \end{bmatrix}$, which
reduces to $\begin{bmatrix} 1 & 0 & 1 \\ 0 & 1 & 1 \\ 0 & 0 & 0 \end{bmatrix}$ ($m=n=3$, $r=2$). The bases:
row space $(1,0,1),(0,1,1)$; nullspace $\svec=\begin{bsmallmatrix} -1\\-1\\1 \end{bsmallmatrix}$; columns
$\aaa_1=\begin{bsmallmatrix} 1\\2\\3 \end{bsmallmatrix},\aaa_2=\begin{bsmallmatrix} 1\\1\\2 \end{bsmallmatrix}$;
left nullspace $\yy=\begin{bsmallmatrix} 1\\1\\-1 \end{bsmallmatrix}$. Verify the right angles with dot products:
\[
  (1,0,1)\cdot\svec=\blank{0.9cm},\quad (0,1,1)\cdot\svec=\blank{0.9cm}
  \qquad\Rightarrow\qquad N(A)\perp C(A^{\T});
\]
\[
  \yy\cdot\aaa_1=\blank{0.9cm},\quad \yy\cdot\aaa_2=\blank{0.9cm}
  \qquad\Rightarrow\qquad N(A^{\T})\perp C(A).
\]
\textbf{Next (Unit 4): Orthogonality.} These right angles are the whole idea behind \emph{projection} and
\emph{least squares} --- finding the closest point in a subspace. The Fundamental Theorem is the bridge from
Unit~3 (the four subspaces) to Unit~4.
\end{notesbox}

\begin{practicebox}
Show your reasoning.
\begin{enumerate}[leftmargin=1.6em, itemsep=2pt]
  \item State both parts of the Fundamental Theorem for a $4\times 6$ matrix of rank $r=3$: the four dimensions, and which two pairs of subspaces are perpendicular (and in which room). \writeline
  \item For $A=\begin{bmatrix} 1 & 2 \\ 2 & 4 \end{bmatrix}$, the nullspace vector is $\svec=\begin{bsmallmatrix} -2\\1 \end{bsmallmatrix}$ and a row is $(1,2)$. Show by a dot product that $\svec\perp$ the row space. \writeline
  \item In one sentence: why does $A\xx=\zero$ force $\xx$ to be perpendicular to every row of $A$? \writeline
\end{enumerate}
\end{practicebox}

\end{document}
